% !TEX root = ../algorithms.tex

\section{Основы линейного программирования}

\newcommand{\exercise}{\par\noindent\textbf{Упражнение.}\par}

\begin{Definition}{Множество допустимых решений}{}
Множество $P = \{x : Ax \le b\}$ называется множеством допустимых решений. 
Данное множество является \textbf{выпуклым}, то есть:
\[
\forall x, y \in P, \forall \lambda \in (0, 1) \implies \lambda x + (1 - \lambda) y \in P
\]
Это следует из того, что $P$ является пересечением полупространств, а полупространства — выпуклые множества.
\end{Definition}

\exercise
Докажите, что пересечение выпуклых множеств выпукло.

\begin{Definition}{Оптимальное решение}{}
Точка $x^* \in \arg\max_{x \in P} c^T x$ называется оптимальным решением задачи линейного программирования.
\end{Definition}

\subsection{Ограниченность и вершины}

Линейная программа называется \textbf{ограниченной}, если $P \neq \varnothing$ и $\exists M : \forall x \in P \implies c^T x \le M$ (в случае минимизации $c^T x \ge M$).

\begin{Definition}{Вершина}{}
Точка $x \in P$ называется вершиной, если $\nexists y, z \in P, y \neq z, \exists \lambda \in (0, 1):$
\[
x = \lambda y + (1 - \lambda) z
\]
\end{Definition}

Интуитивно понятно, что если в задаче линейного программирования есть оптимальное решение, то существует оптимальное решение, являющееся вершиной.

\begin{Theorem}{}{}
Можно доказать, что если линейная программа ограничена...
\end{Theorem}

\section{Доказательство леммы о разделяющей гиперплоскости}

\begin{Theorem}{}{}
Весь конус $K$ лежит по другую сторону от гиперплоскости, т.е. $Az \le 0$, в то время как $c^T z > 0$.
\end{Theorem}

\textbf{Доказательство:}

I случай: $c \notin K$ (т.е. условие 1 не выполняется).
Пусть $z = c - \hat{y}$, где $\hat{y}$ — это ближайшая к $c$ точка из $K$. 
Как построить $\hat{y}$? 
$\forall y \in K$ выберем $y_k \in K$ такие, что $|c - y_k| \le \inf |c - y| + \frac{1}{k}$. 
По теореме Больцано-Вейерштрасса выберем подпоследовательность $\{y_{n_k}\} \to \hat{y}$. Можно доказать, что $K$ замкнут, поэтому $\hat{y} \in K$.

Рассмотрим подпространство $L = \operatorname{span}(c, \hat{y})$, изоморфное $\RR^2$. В силу школьной геометрии в $\RR^2$, вектор $z = c - \hat{y}$ перпендикулярен касательной к конусу в точке $\hat{y}$. 

II случай: Если подпространство $L$ изоморфно $\RR^3$.
Предположим, $c^T z > 0$. Если конус $K$ — выпуклый, то $K$ содержит отрезок
\[
[a_i, \hat{y}] = \{ s a_i + (1-s)\hat{y} \mid s \in [0, 1] \}.
\]
Так как $\hat{y}$ — ближайшая, с положительной стороны плоскости $\{x : x^T z = 0\}$ точек быть не может. Из элементарной геометрии в $\RR^3$ следует противоречие.

\exercise
Из геометрии в $\RR^3$ следует противоречие с тем, что $\hat{y}$ — ближайшая точка.

\section{Теорема двойственности}

Для $A \in \RR^{m \times n}$, $b \in \RR^m$, $c \in \RR^n$ рассмотрим прямую и двойственную задачи линейного программирования:

\begin{equation*}
    \begin{array}{c|c}
        \text{Прямая задача} & \text{Двойственная задача} \\ \hline
        \max c^T x & \min b^T y \\
        Ax \le b & A^T y = c \\
        & y \ge 0
    \end{array}
\end{equation*}

\begin{Theorem}{Двойственности}{}
Если одна из задач имеет оптимальное решение, то и другая имеет решение, причем $c^T x^* = b^T y^*$, где $x^*, y^*$ — оптимальные решения.
\end{Theorem}

\begin{Corollary}{Следствие}{}
Если целевая функция одной из задач не ограничена, то другая задача не имеет допустимых решений.
\end{Corollary}

\subsection{Доказательство}

1) Пусть $x, y$ — допустимые решения. Тогда:
\[
c^T x = (A^T y)^T x = y^T (Ax) \le y^T b = b^T y
\]
Мы получили \textbf{слабую теорему двойственности}: значение целевой функции прямой задачи на любом допустимом решении не превосходит значения двойственной задачи.

Из этого замечаем, что если обе задачи имеют допустимые решения, то они ограничены и имеют оптимальные решения.

\section{Доказательство сильной теоремы двойственности}

Пусть $x^*$ — оптимальное решение прямой задачи $Ax \le b, \max c^T x$. 
Перестановкой строк в системе ограничений добьемся того, что
\[
A = \begin{pmatrix} A' \\ A'' \end{pmatrix}
\quad \text{и} \quad
b = \begin{pmatrix} b' \\ b'' \end{pmatrix},
\]
где:
\[
A' x^* = b', \qquad A'' x^* < b''.
\]

Заметим, что $\nexists z : A' z \le 0$ и $c^T z > 0$ (т.к. иначе при малом $\eps > 0$ точка $x^{**} = x^* + \eps z$ удовлетворяла бы условию $A x^{**} \le b$, но $c^T x^{**} > c^T x^*$, что противоречит оптимальности $x^*$).

По лемме Фаркаша $\exists y' \ge 0 : (A')^T y' = c$. 
Пусть
\[
y = \begin{pmatrix} y' \\ 0 \end{pmatrix}.
\]
Тогда $y \ge 0$ и $A^T y = (A')^T y' + (A'')^T \cdot 0 = c$, то есть $y$ — допустимое решение двойственной задачи.

Проверим значения целевых функций:
\[
b^T y = (b')^T y' = (A' x^*)^T y' = (x^*)^T (A')^T y' = (x^*)^T c = c^T x^*.
\]
Следовательно, $y$ — оптимальное решение двойственной задачи.

\subsection{Симметричность двойственности}

Осталось доказать, что если двойственная задача имеет оптимальное решение, то и прямая тоже. 
Построим двойственную задачу к двойственной. Представим двойственную задачу в каноническом виде:
\begin{gather*}
    \min b^T y \iff \max (-b^T y) \\
    \begin{cases} A^T y = c \\ y \ge 0 \end{cases}
    \iff
    \begin{cases} A^T y \le c \\ -A^T y \le -c \\ -y \le 0 \end{cases}
\end{gather*}
Двойственная к этой задаче будет в точности эквивалентна исходной задаче $\max c^T x, Ax \le b$.

\exercise
Докажите эквивалентность последней конструкции самостоятельно.